\section{Jandex: ler anotações sem carregar classes}
\label{sec:jandex}

Para descobrir beans, endpoints REST e entidades sem carregar classes, o
Quarkus usa o \termo{Jandex}{Biblioteca da SmallRye que indexa arquivos
\texttt{.class}; o índice pode ser gravado em
\texttt{META-INF/jandex.idx} dentro do JAR.}. O repositório do projeto
(\citetitle{doc:jandex}) o descreve como \enquote{a space efficient Java
class file indexer and offline reflection library}: \emph{offline
reflection} porque responde às perguntas que a reflexão responderia
(que classes têm esta anotação, quem implementa esta interface) sem
pedir nada à JVM.

\paragraph{Por que isso é possível.}
Anotações com retenção \texttt{RUNTIME} ficam gravadas no próprio
\texttt{.class}, no atributo \texttt{RuntimeVisibleAnnotations} (JVMS
§4.7.16, \citetitle{doc:jvms25}), e apontam para entradas do constant
pool (seção~\ref{sec:constant-pool}). Um indexador só precisa ler bytes
e montar tabelas; não precisa de class loader, verificação nem
\texttt{<clinit>}. Com \texttt{javap -v} você vê esse atributo em
qualquer classe anotada.

\subsection*{Como uma dependência entra no índice}

O guia \citetitle{doc:quarkus-cdi-reference} descreve uma
\emph{descoberta de beans simplificada}, com um único \emph{bean archive}
formado, entre outras fontes, pela aplicação e por:
\begin{itemize}
\item dependências que têm um descritor \texttt{beans.xml} (o conteúdo é
  ignorado);
\item dependências que já trazem um índice Jandex
  (\texttt{META-INF/jandex.idx});
\item dependências listadas em \texttt{quarkus.index-dependency} no
  \texttt{application.properties}.
\end{itemize}

O guia de extensões completa: o índice das classes analisadas
\enquote{will not automatically include your external class dependencies}
e nenhum artefato é incluído transitivamente. O sintoma clássico é um
módulo compartilhado (por exemplo, adaptadores comuns a vários serviços)
num JAR separado cujos beans \enquote{somem}: o Quarkus não os viu porque
o JAR não estava no índice.

\subsection*{Remoção de beans não usados}

Com o índice e o grafo de injeção em mãos, o container CDI do Quarkus
(ArC) faz algo que só é seguro em build time. Segundo o guia de CDI, ele
tenta remover beans, interceptors e decorators não usados durante o
build, por padrão. Dois avisos do mesmo guia:
\begin{itemize}
\item O Quarkus não detecta busca programática via \texttt{CDI.current()};
  um bean usado só assim pode ser removido. Desligue com
  \texttt{quarkus.arc.remove-unused-beans=none} ou, melhor, injete o bean.
\item Classes sem \emph{bean defining annotation} (por exemplo,
  \texttt{@ApplicationScoped}) não são descobertas. Isso importa na
  arquitetura do capítulo~\ref{cap:arquitetura-quarkus}: uma classe de
  caso de uso sem anotação precisa de um \emph{producer} para virar bean.
\end{itemize}

\paragraph{O preço.}
O guia de CDI admite o custo dessa abordagem: o ArC não implementa o
\emph{CDI Full}, e \emph{CDI Portable Extensions} não são suportadas;
o guia apresenta isso como a desvantagem das otimizações de build time.
Biblioteca que
funciona via portable extension em outro servidor precisa de uma
extensão Quarkus equivalente.

\paragraph{Vendo o resultado.}
Em dev mode, os endpoints \texttt{/q/arc} e \texttt{/q/arc/beans} mostram
os beans que sobraram, e o nível de log \texttt{DEBUG} mostra os que foram
removidos, segundo o mesmo guia.
